\documentclass[12pt]{article}
\usepackage[utf8]{inputenc}
\usepackage[T1]{fontenc}
\usepackage{hyperref}
\title{Zimba Language Resources for K-12 Curriculum Development}
\author{Compiled by Claude Sonnet 5.5 (claude-sonnet-5-5) for Andr\'as Kornai}
\date{}
\begin{document}
\maketitle

\section{Introduction}
Zimba (ISO 639-3 zmb; Glottocode zimb1251) is a Bantu language of the Democratic Republic of the Congo, spoken south of the Elila River. Ethnologue places it in the Lega-Holoholo (D.26) subgroup. Wikipedia cites about 120,000 speakers. Alternate names include Binja, Kibinja, KiZimba, Kyenyemamba and Mamba, and dialects named include Kwange, Semulu and South Binja. It should not be confused with Zemba (Dhimba) of Angola. The searches found little material written in or for Zimba, and nothing designed for schools.

\section{Primary Educational Resources}
No primers, textbooks or school programmes for Zimba turned up in the searches. The language is not known to be taught in schools. A curriculum would have to build its own materials from community and field sources.

\section{Linguistic and Reference Materials}
\begin{itemize}
\item Wikipedia, Zimba language (overview, classification, dialects): \url{https://en.wikipedia.org/wiki/Zimba_language}
\item Wikipedia, Lega--Binja languages (genetic context): \url{https://en.wikipedia.org/wiki/Lega%E2%80%93Binja_languages}
\item Ethnologue, Lega-Holoholo subgroup listing: \url{https://ethnologue.com/subgroup/4191/}
\item ASJP wordlist entry for Zimba: \url{https://asjp.clld.org/languages/ZIMBA}
\end{itemize}

\section{Religious and Community Resources}
\begin{itemize}
\item Joshua Project, Zimba people group in DR Congo (names, dialects, Scripture status; it reports a Bible portion dated 2022 and no audio Bible): \url{https://joshuaproject.net/people_groups/16105/CG}
\item Joshua Project, language page for zmb: \url{https://joshuaproject.net/index.php/languages/zmb}
\end{itemize}

\section{Digital and Audio Resources}
\begin{itemize}
\item Global Recordings Network language page (the search summary reported no recordings in Zimba): \url{https://glorecnet-website.k.globalrecordings.net/en/language/18717}
\end{itemize}

\section{Curriculum Development Context}
The sources found are overview and database entries only. Anyone planning a Zimba curriculum would first need to work with the Zimba community and speakers to document the language, agree an orthography, and develop materials. French is the likely language of instruction in DR Congo, and regional Bantu lingua francas are also in use. The searches did not confirm any Zimba language policy or school programme.

\end{document}